\documentclass[11pt,a4paper]{article}

\usepackage[margin=1in]{geometry}
\usepackage[utf8]{inputenc}
\usepackage[T1]{fontenc}
\usepackage{hyperref}
\usepackage{enumitem}
\usepackage{graphicx}
\usepackage{xcolor}
\usepackage{titlesec}
\usepackage{array}
\usepackage{colortbl}
\usepackage{longtable}
\usepackage{fancyhdr}
\usepackage{parskip}
\usepackage{needspace}
\usepackage{etoolbox}

% -----------------------------------------------------
% Formatting
% -----------------------------------------------------

\definecolor{rowshade}{RGB}{235,239,245}

\titleformat{\section}
  {\normalfont\Large\bfseries\color{black}}
  {\thesection.}{0.6em}{}
  [{\titlerule[1pt]}]

\titleformat{\subsection}
  {\normalfont\large\bfseries\color{black}}
  {\thesubsection}{0.6em}{}

\titlespacing*{\section}
  {0pt}{2.2ex plus 1ex minus .2ex}{1.2ex}

\titlespacing*{\subsection}
  {0pt}{1.6ex plus .5ex minus .2ex}{0.8ex}

% Keep a heading together with the text that follows it.
% Force a page break before the heading rather than leaving it
% orphaned at the bottom of a page.
\pretocmd{\section}{\needspace{7\baselineskip}}{}{}
\pretocmd{\subsection}{\needspace{4\baselineskip}}{}{}

\setlist[itemize]{
    leftmargin=1.4em,
    itemsep=1pt,
    topsep=2pt,
    parsep=0pt
}

\setlist[enumerate]{
    leftmargin=1.4em,
    itemsep=1pt,
    topsep=2pt,
    parsep=0pt
}

\setlength{\parskip}{3pt}

\hypersetup{
    colorlinks=false,
    hidelinks
}

% Header / footer
\pagestyle{fancy}
\fancyhf{}
\renewcommand{\headrulewidth}{0.4pt}
\fancyhead[L]{\small Intelligent Expense Tracker - Project Proposal}
\fancyhead[R]{\small\thepage}
\setlength{\headheight}{14pt}

% -----------------------------------------------------

\begin{document}

% =====================================================
% TITLE
% =====================================================

\begin{center}

    {\huge\bfseries
    Intelligent Expense Tracker}\\[0.5em]

    {\large
    An Automated Platform for Expense Capture, Collaborative Spending,
    and Financial Intelligence}
    \\[0.9em]

    {\normalsize\itshape
    Software Engineering - Semester Project Proposal}\\[0.8em]

    {\small\bfseries Team Members}\\[0.3em]

    {\small
    Aashi Tyagi (1024170263) \quad
    Sarwangini Hamal (1024170262)\\[0.15em]
    Divyansh Kumar (1024170267) \quad
    Gauransh Arora (1024170264)}\\[0.5em]

    {\small\textbf{Submitted to:} Nisha Thakur}

\end{center}

\vspace{0.8em}

% =====================================================
% 1. INTRODUCTION
% =====================================================

\section{Introduction}

Every digital payment already carries the information needed to
understand it: a UPI transaction has a merchant, an amount, and a time;
a printed bill lists individual items and prices; a bank statement
records a running history of debits and credits. In practice, however,
this information stays scattered across screenshots, paper receipts, and
PDF statements, and none of it is usable until a person manually
retypes it into an expense-tracking app.

This project proposes a system built around a simple idea: instead of
asking users to record their expenses, the application should read the
financial documents users already generate and turn them into a single,
structured financial history. Expense splitting, analytics, and
financial insights are then built on top of this history rather than on
manually typed entries. The project is therefore not simply an expense
tracker or a bill-splitting tool, but an automated financial record
and analysis system, of which splitting and analytics are two of
several components.

% =====================================================
% 2. PROBLEM STATEMENT
% =====================================================

\section{Problem Statement}

Financial information is generated constantly, but it is not structured
or easily usable. This produces several concrete problems:

\begin{itemize}

    \item \textbf{Fragmented financial information.} Receipts contain
    merchant, item, and price detail; UPI screenshots contain
    transaction detail; bank statements contain a running transaction
    history. None of these sources are connected, so a user's complete
    spending picture never exists in one place.

    \item \textbf{Repeated manual effort.} Even though this information
    already exists digitally, users must retype it into a tracker,
    which is slow enough that most people stop doing it consistently.

    \item \textbf{Limited understanding of spending.} Most expense
    trackers can answer ``how much did I spend'', but not ``where did
    it go'', ``why did it increase'', or ``what is recurring''.

    \item \textbf{No decision support.} Recorded expenses are rarely
    used to evaluate a planned purchase, and shared expenses add a
    further layer of complexity because the person who pays and the
    people who benefit are often different.

\end{itemize}

The proposed system addresses this by automating the capture step and
then using the resulting structured history to answer these harder
questions, rather than stopping at transaction logging.

% =====================================================
% 3. OBJECTIVES
% =====================================================

\section{Objectives}

\noindent\textbf{Primary goal.}
To design and build a platform that automatically converts bills, UPI
screenshots, and bank statements into a structured, verified financial
history, and uses that history for collaborative expense management and
spending intelligence.

\noindent\textbf{Specific objectives:}

\begin{itemize}

    \item Automate expense capture from bills, receipts, UPI
    screenshots, and bank statements using OCR.

    \item Convert extracted information into structured, categorized
    transactions, including merchant, amount, date, and individual
    items, with user verification before storage.

    \item Support collaborative expense management through equal,
    percentage, and custom splitting, with debt tracking between users.

    \item Provide multi-dimensional spending analytics across
    categories, merchants, items, and time periods.

    \item Detect recurring expenses and unusual spending patterns from
    transaction history.

    \item Support spending decisions through purchase-readiness
    analysis and wishlist tracking.

    \item Enable natural-language interaction with the user's own
    financial history.

\end{itemize}

% =====================================================
% 4. PROPOSED SOLUTION
% =====================================================

\section{Proposed Solution}

At a high level, the system takes any financial document a user already
has, such as a bill, a UPI payment screenshot, or a bank statement, and
turns it into a verified transaction. The user shares or uploads the
document; the system extracts and normalizes the relevant fields and
suggests a category; the user confirms or corrects this before it is
saved. Once confirmed, the transaction becomes part of the user's
financial history, which is then used by separate modules for
collaborative splitting, analytics, recurring and unusual spending
detection, purchase-readiness analysis, and natural-language queries.

\begin{figure}[ht]
    \centering
    \includegraphics[width=0.75\textwidth]{system_flow}
    \caption{Overall System Flow}
    \label{fig:system-flow}
\end{figure}

This flow keeps the user in control of what enters their financial
record, while shifting the repetitive work of data entry onto the
system.


% =====================================================
% 5. SYSTEM ARCHITECTURE AND METHODOLOGY
% =====================================================

\section{System Architecture and Methodology}

\subsection{System Architecture}

The system follows a three-tier architecture: a Flutter mobile frontend,
a FastAPI backend exposing REST APIs, and a PostgreSQL database for
persistent storage.

\begin{figure}[ht]
    \centering
    \includegraphics[width=0.75\textwidth]{backend_architecture}
    \caption{System Architecture}
    \label{fig:backend-architecture}
\end{figure}

The backend is organized into an extraction layer for OCR and document
parsing, a processing layer for categorization and verification, and an
intelligence layer for splitting, analytics, pattern detection, and
queries. These components read from and write to a shared PostgreSQL
schema covering users, transactions, items, categories, splits, and
wishlist entries.

\subsection{Automated Expense Extraction and Verification}

The core extraction workflow is:

\begin{enumerate}

    \item A user uploads a bill or shares a UPI payment screenshot.

    \item OCR extracts raw text from the image.

    \item The processing layer identifies merchant, amount, date, items,
    quantities, and taxes from the extracted text.

    \item Transactions and individual items are categorized using
    merchant information, keywords, and lightweight classification.

    \item The extracted information, along with a confidence score, is
    shown to the user for verification.

    \item The confirmed transaction is stored in PostgreSQL and becomes
    available to every downstream module.

\end{enumerate}

The system does not treat OCR output as final. Extracted data is always
presented for user confirmation before it becomes part of the financial
history, which limits how far an extraction error can propagate.

\subsection{Collaborative Expense Management}

Shared expenses such as meals, trips, hostel costs, or group purchases
can be split between participants using equal, percentage-based, or custom
amounts. The system tracks who paid, who owes what, and maintains a
running settlement history between users, so debts do not need to be
tracked manually outside the app.

\subsection{Expense Tracking and Spending Analytics}

The application maintains a chronological expense timeline containing
merchant, amount, category, items, and the originating bill or
screenshot. The analytics layer builds on this timeline to provide
category-wise and monthly spending, merchant-level and item-level
analysis, comparisons with previous months, and ``where did my money
go'' summaries.

% =====================================================
% 6. FINANCIAL INTELLIGENCE
% =====================================================

\section{Financial Intelligence}

This section covers the modules that use the stored transaction history
to go beyond recording expenses and actually derive information from
them. These modules are the main differentiator of the project over a
standard tracker or splitting app.

\subsection{Unusual Spending Detection}

Historical spending per category is used to establish a normal range for
the user. Transactions or category totals that deviate noticeably from
this range are flagged for the user's attention. This is a
historical-pattern-based check rather than a trained anomaly-detection
model.

\subsection{Recurring Expense and Subscription Detection}

Transactions with similar merchants, amounts, and regular time intervals
are identified as likely recurring expenses or subscriptions. This is
inferred purely from transaction history rather than from direct access
to UPI mandates or bank-level subscription data.

\subsection{Purchase Readiness and Wishlist Analysis}

Users can add a planned purchase to a wishlist along with its expected
price. The system evaluates the purchase against current spending,
monthly average spending, recurring payments, historical trends, savings
goals, and discretionary budget, and returns a simple verdict:
\textbf{Comfortable}, \textbf{Possible}, or
\textbf{Not recommended currently}.

The system can also compare current-month spending with the user's
historical average to identify whether the user currently has additional
discretionary room for a wishlist purchase. This is intended as a
data-driven spending indicator, not professional financial advice.

\subsection{Ask Your Expenses}

An ``Ask Your Expenses'' interface lets users query their own financial
history in natural language instead of navigating separate dashboards.

Example questions include:

\begin{itemize}
    \item ``How much did I spend on food this month?''
    \item ``What are my recurring monthly expenses?''
    \item ``How much does Rahul owe me?''
    \item ``Where did most of my money go this month?''
\end{itemize}

Supported questions are mapped to structured database queries or
analytical operations and executed only against the user's own data.

\subsection{Bank Statement Processing}

Supported PDF or CSV bank statements can be uploaded and converted into
structured transactions, which are normalized, categorized, and merged
with existing expenses after user verification. Real-time bank
synchronization is outside the initial scope, since it requires
additional financial ecosystem access and consent mechanisms.

% =====================================================
% 7. SCOPE AND MVP
% =====================================================

\section{Scope and MVP}

\subsection{MVP Scope}

The MVP includes bill and UPI screenshot extraction with user
verification, expense categorization and history, collaborative
splitting, spending analytics, recurring-expense detection, unusual
spending detection, wishlist and purchase-readiness analysis, the
Ask Your Expenses query interface, and supported bank-statement
processing.

\subsection{Out of Scope}

Real-time bank and UPI account synchronization, direct control or
blocking of third-party payments, and professional financial advisory or
investment-recommendation features are outside the MVP. These require
regulatory and financial ecosystem access beyond a semester project and
are noted as future extensions instead.

% =====================================================
% 8. TESTING AND EVALUATION
% =====================================================

\section{Testing and Evaluation}

Success will be measured against four concrete evaluation areas, chosen
to match the project's central claims:

\begin{itemize}

    \item \textbf{OCR extraction accuracy.} Extracted merchant, amount,
    date, and item fields will be compared against manually verified
    ground truth on a sample set of bills and screenshots.

    \item \textbf{Categorization accuracy.} The proportion of
    transactions assigned to the correct category, checked against
    user corrections.

    \item \textbf{Time saved.} The time taken to record an expense via
    upload and verification will be compared against fully manual entry.

    \item \textbf{Query correctness.} A fixed set of natural-language
    questions will be run against known data, and returned answers
    checked against the underlying database calculations.

\end{itemize}

Unit tests will cover backend services and transaction-processing logic,
and integration tests will verify the complete document-to-database
workflow before user testing begins.

% =====================================================
% 9. TECHNOLOGY STACK AND RESOURCES
% =====================================================

\section{Technology Stack and Resources}

\noindent\textbf{Technology stack.}
Flutter will provide the cross-platform mobile interface, covering
authentication, screenshot sharing, image upload, expense verification,
dashboards, and financial queries. FastAPI will expose REST APIs for
authentication, transaction processing, splitting, analytics, and
financial queries. PostgreSQL will store users, transactions, items,
categories, splits, and wishlist data. Git and GitHub will be used for
version control.

\noindent\textbf{Human resources.}
The team will divide responsibilities across mobile development, backend
development, database design, OCR and data processing, and testing.

\noindent\textbf{Budget.}
No external funding is required for the prototype. OCR, PostgreSQL, and
the backend can be implemented using free or open-source resources, with
free-tier deployment for the staging build. Paid AI services and
production financial integrations remain optional future requirements.

% =====================================================
% 10. TIMELINE
% =====================================================

\section{Timeline}

The project will be developed incrementally, with the automated capture
workflow completed before collaborative and intelligence features.

{\small
\renewcommand{\arraystretch}{1.15}
\setlength{\tabcolsep}{5pt}

\begin{center}

\begin{tabular}{
    @{}
    >{\raggedright\arraybackslash}p{1.5cm}
    >{\raggedright\arraybackslash}p{9.3cm}
    >{\raggedright\arraybackslash}p{3.6cm}
    @{}
}

\hline

\rowcolor{rowshade}
\textbf{Weeks} & \textbf{Phase} & \textbf{Deliverable} \\

\hline

1--2 &
Planning and Design: requirements, architecture, UI &
ER diagram \\

\hline

3--4 &
Core Application: authentication, transaction model, manual entry &
Mobile shell + backend \\

\hline

5--6 &
Automated Extraction: OCR, bill and UPI screenshot processing &
Expense extraction module \\

\hline

7--8 &
Expense Management: categories, splitting, transaction history &
Core expense tracker \\

\hline

9--10 &
Analytics: dashboards, category, merchant, and item spending &
Analytics dashboard \\

\hline

11--12 &
Financial Intelligence: recurring and unusual spending detection,
wishlist, and queries &
Intelligence modules \\

\hline

13 &
Testing and Integration: unit, integration, and user testing &
Integrated build \\

\hline

14 &
\textbf{Buffer}: bug fixing and optimization &
Stabilised build \\

\hline

15 &
Deployment, documentation, and demonstration &
Final application + report \\

\hline

\end{tabular}

\end{center}
}

% =====================================================
% 11. RISKS AND FUTURE EXTENSIONS
% =====================================================

\section{Risks and Future Extensions}

\subsection{Risks and Mitigation}

\begin{itemize}

    \item \textbf{OCR inaccuracies}, especially for low-quality or
    unusual receipt formats, are mitigated through confidence scores and
    mandatory user verification before storage.

    \item \textbf{Inconsistent receipt and screenshot layouts} across
    merchants are mitigated by supporting common formats first and
    providing a manual-correction fallback.

    \item \textbf{Incorrect categorization} affecting analytics is
    mitigated through user correction, which is fed back into
    classification rules.

    \item \textbf{Financial data sensitivity} is mitigated through
    authentication, access control, secure transport, and data
    minimisation.

    \item \textbf{Scope creep}, given the number of modules, is
    mitigated through an MVP-first approach with clearly defined
    out-of-scope features.

\end{itemize}

\subsection{Future Extensions}

Future directions include consent-based bank and UPI account aggregation,
more advanced anomaly detection, household or shared-family expense
management, and more sophisticated budgeting and financial
recommendations, once sufficient real usage data exists to support them
responsibly.

A future parent dashboard could allow parents to define spending limits,
monitor a child's spending, and receive alerts when a child approaches a
budget threshold. However, directly blocking payments through
third-party UPI applications is outside the core scope.

% =====================================================
% 12. SUMMARY
% =====================================================

\section{Summary}

This project proposes the Intelligent Expense Tracker, a system that
automates the capture of financial information from bills, UPI
screenshots, and bank statements, and turns it into a verified,
structured financial history. The central objective is to transform
fragmented and unstructured transaction information into a structured
financial history that can be automatically analyzed to help users
understand and make better everyday spending decisions through
collaborative splitting, spending analytics, recurring and unusual
spending detection, purchase-readiness analysis, and natural-language
queries, built on Flutter, FastAPI, and PostgreSQL.

\end{document}